%!TEX program = lualatex
% SPDX-License-Identifier: CC-BY-SA-4.0
% A maximalist, entirely vector specimen. ./build-beamer.sh --grand
\documentclass[aspectratio=169,11pt]{beamer}
\usetheme{AtlasGrand}
\usepackage{booktabs}
\usepackage{array}
% SPDX-License-Identifier: MIT
% Copyright (c) 2026 the repository contributors
% Engraved mathematical plates for the grand specimen. Entirely vector.
% The seven-edge tree and six states agree with tree-pursuit.tex.
\newcommand{\GrandTree}[1]{%
  \begingroup
  \ifcase#1\relax
  \or\def\gc{a}\def\gr{d}\def\region{b,c,d,e,f,g,h}\def\edges{a/b,b/c,c/d,d/e,b/f,b/g,c/h}%
  \or\def\gc{b}\def\gr{d}\def\region{c,d,e,h}\def\edges{b/c,c/d,d/e,c/h}%
  \or\def\gc{b}\def\gr{e}\def\region{c,d,e,h}\def\edges{b/c,c/d,d/e,c/h}%
  \or\def\gc{c}\def\gr{e}\def\region{d,e}\def\edges{c/d,d/e}%
  \or\def\gc{d}\def\gr{e}\def\region{e}\def\edges{d/e}%
  \or\def\gc{e}\def\gr{e}\def\region{}\def\edges{}%
  \fi
  \begin{tikzpicture}[x=1cm,y=1cm]
    \path[use as bounding box] (-.65,-2.05) rectangle (5.8,2.35);
    \begin{scope}[overlay]
      \fill[AtlasPaper!85!white] (-.65,-2.05) rectangle (5.8,2.35);
      \draw[AtlasGold,line width=.7pt] (-.65,-2.05) rectangle (5.8,2.35);
      \draw[AtlasGold,line width=.2pt] (-.57,-1.97) rectangle (5.72,2.27);
      \foreach \x in {0,.4,...,5.4}
        \draw[AtlasGold,opacity=.13,line width=.2pt] (\x,-1.75)--(\x,1.9);
      \foreach \y in {-1.6,-1.2,...,1.6}
        \draw[AtlasGold,opacity=.13,line width=.2pt] (-.4,\y)--(5.55,\y);
      \node[fill=AtlasPaper!85!white,text=AtlasRubric,font=\scriptsize\scshape]
        at (2.575,2.27) {Tabula I\quad Arbor finita};
      \node[fill=AtlasPaper!85!white,text=AtlasGold,font=\scriptsize\itshape]
        at (2.575,-1.97) {The branches remember every advance};
      \coordinate (a) at (0,0);
      \coordinate (b) at (1.25,0);
      \coordinate (c) at (2.5,0);
      \coordinate (d) at (3.75,.9);
      \coordinate (e) at (5,.9);
      \coordinate (f) at (1.25,1.35);
      \coordinate (g) at (1.25,-.95);
      \coordinate (h) at (2.5,-.95);
      \foreach \u/\v in {a/b,b/c,c/d,d/e,b/f,b/g,c/h}
        \draw[AtlasInk!28!AtlasPaper,dashed,line width=.55pt] (\u)--(\v);
      \foreach \u/\v in \edges {
        \draw[AtlasGold!45!AtlasPaper,line width=2.8pt] (\u)--(\v);
        \draw[AtlasInk,line width=.55pt,double=AtlasPaper,double distance=.8pt] (\u)--(\v);
      }
      \foreach \v in {a,b,c,d,e,f,g,h}
        \fill[AtlasInk!35!AtlasPaper] (\v) circle (.035);
      \foreach \v in \region {
        \fill[AtlasGold!16!AtlasPaper] (\v) circle (.24);
        \draw[AtlasGold,line width=.35pt] (\v) circle (.24);
        \draw[AtlasGold!60!AtlasPaper,line width=.2pt] (\v) circle (.29);
        \fill[AtlasInk] (\v) circle (.055);
      }
      \foreach \v in {a,d,e,g,h}
        \node[below=3.7mm,text=AtlasInk,font=\small\itshape] at (\v) {\v};
      \foreach \v in {b,c}
        \node[below right=3mm,text=AtlasInk,font=\small\itshape] at (\v) {\v};
      \node[above=3mm,text=AtlasInk,font=\small\itshape] at (f) {f};
      \ifdefstrequal{\gc}{\gr}{%
        \node[circle,draw=AtlasGold,fill=AtlasRubric,text=AtlasPaper,
          double=AtlasPaper,double distance=1pt,minimum size=8mm,
          inner sep=0pt,font=\scriptsize] at (\gc) {C/R};
      }{%
        \node[circle,draw=AtlasGold,fill=AtlasRubric,text=AtlasPaper,
          double=AtlasPaper,double distance=1pt,minimum size=7mm,
          inner sep=0pt,font=\small] at (\gc) {C};
        \node[diamond,draw=AtlasInk,fill=AtlasPaper,text=AtlasInk,
          double=AtlasGold,double distance=.7pt,minimum size=9mm,
          inner sep=0pt,font=\small] at (\gr) {R};
      }
    \end{scope}
  \end{tikzpicture}%
  \endgroup
}

\newcommand{\GrandDial}[1]{%
  \begin{tikzpicture}[baseline=(number.base)]
    \path[use as bounding box] (-1.05,-1.05) rectangle (1.05,1.05);
    \foreach \i in {1,...,7}{
      \pgfmathsetmacro{\start}{90-(\i-1)*360/7-3}
      \pgfmathsetmacro{\stop}{90-\i*360/7+3}
      \ifnum\i>#1\relax\def\sector{AtlasGold!15!AtlasPaper}
      \else\def\sector{AtlasRubric}\fi
      \fill[\sector] (\start:.9) arc (\start:\stop:.9)--(\stop:.76)
        arc (\stop:\start:.76)--cycle;
    }
    \draw[AtlasGold,line width=.3pt] (0,0) circle (.96);
    \draw[AtlasGold,line width=.3pt] (0,0) circle (.69);
    \node[text=AtlasRubric,font=\fontsize{27}{27}\selectfont] (number) at (0,.06) {#1};
    \node[text=AtlasGold,font=\tiny\scshape] at (0,-.4) {remaining};
  \end{tikzpicture}%
}

\newcommand{\GrandCycle}[3]{%
  % Cop angle, robber angle, second cop angle (-1 = absent).
  \begin{tikzpicture}[x=1cm,y=1cm]
    \path[use as bounding box] (-2.95,-2.75) rectangle (2.95,2.75);
    \begin{scope}[overlay]
      \draw[AtlasGold,line width=.4pt] (0,0) circle (2.3);
      \draw[AtlasGold,line width=.2pt] (0,0) circle (2.38);
      \foreach \a in {0,10,...,350}
        \draw[AtlasGold,line width=.3pt] (\a:2.3)--(\a:2.36);
      \foreach \a in {0,60,...,300}{
        \pgfmathsetmacro{\b}{\a+60}
        \draw[AtlasGold,line width=2.5pt] (\a:1.75)--(\b:1.75);
        \draw[AtlasInk,line width=.55pt] (\a:1.75)--(\b:1.75);
        \node[circle,fill=AtlasInk,inner sep=1.8pt] at (\a:1.75) {};
        \node[text=AtlasGold,font=\scriptsize] at (\a:2.6) {\a};
      }
      \node[text=AtlasGold,font=\Large] at (0,.15) {$C_6$};
      \node[text=AtlasGold,font=\tiny\scshape] at (0,-.3) {circuitus};
      \ifnum#3>-1\relax
        \node[circle,fill=AtlasRubric,draw=AtlasGold,text=AtlasPaper,
          double=AtlasPaper,double distance=.7pt,minimum size=7mm,
          inner sep=0pt] at (#3:1.75) {C};
      \fi
      \ifnum#1=#2\relax
        \node[circle,fill=AtlasRubric,draw=AtlasGold,text=AtlasPaper,
          minimum size=8mm,font=\scriptsize,inner sep=0pt] at (#1:1.75) {C/R};
      \else
        \node[circle,fill=AtlasRubric,draw=AtlasGold,text=AtlasPaper,
          double=AtlasPaper,double distance=.7pt,minimum size=7mm,
          inner sep=0pt] at (#1:1.75) {C};
        \node[diamond,fill=AtlasPaper,draw=AtlasInk,text=AtlasInk,
          double=AtlasGold,double distance=.7pt,minimum size=9mm,
          inner sep=0pt] at (#2:1.75) {R};
      \fi
    \end{scope}
  \end{tikzpicture}%
}

\newcommand{\GrandDescent}{%
  \begin{tikzpicture}[x=1.28cm,y=.34cm]
    \path[use as bounding box] (-.4,-3) rectangle (7.6,8.4);
    \foreach \y in {0,...,7}{
      \draw[AtlasGold!28!AtlasPaper,line width=.3pt] (0,\y)--(7.2,\y);
      \node[left,text=AtlasGold,font=\scriptsize] at (0,\y) {\y};
    }
    \foreach \x/\h/\label in {.5/7/a,2/4/b,3.5/2/c,5/1/d,6.5/0/e}{
      \ifnum\h>0
        \fill[AtlasGold!14!AtlasPaper] (\x,0) rectangle ++(.65,\h);
        \fill[pattern=north east lines,pattern color=AtlasGold] (\x,0) rectangle ++(.65,\h);
        \draw[AtlasInk,line width=.5pt] (\x,0) rectangle ++(.65,\h);
      \else
        \draw[AtlasRubric,line width=1.2pt] (\x,0)--++(.65,0);
      \fi
      \node[above,text=AtlasRubric,font=\large] at (\x+.325,\h) {\h};
      \node[below=3mm,text=AtlasInk,font=\small\itshape] at (\x+.325,0) {\label};
    }
    \draw[AtlasRubric,line width=.75pt] (.825,7)--(2.325,4)--(3.825,2)--(5.325,1)--(6.825,0);
    \node[text=AtlasGold,font=\scriptsize\scshape] at (3.6,-2.55) {position of the advancing cop};
  \end{tikzpicture}%
}

\title[Atlas of Pursuit]{Atlas of Pursuit}
\subtitle{A grand discourse on trees, territories, and inevitable capture}
\author{latex-atlas}
\institute{The mathematical cabinet}
\date{4 November 2026}

\newcommand{\GrandBeat}[3]{%
  {\color{AtlasRubric}\scshape\strut #1\strut\par}\smallskip
  \parbox[t][1.8cm][t]{\linewidth}{\raggedright\strut #2\strut}\par
  \smallskip\centerline{\GrandDial{#3}}%
}

\begin{document}

% The frontispiece has its own architectural composition.
{\AtlasGrandNight
\begin{frame}[plain,noframenumbering]
  \begin{tikzpicture}[remember picture,overlay]
    \node[anchor=north west,text=AtlasGoldLight,font=\small]
      at ([xshift=16mm,yshift=-15mm]current page.north west)
      {\textls[160]{\textsc{De venatione in arboribus}}};
    \node[anchor=north west,text=AtlasIvory,inner sep=0pt,align=left]
      at ([xshift=17mm,yshift=-27mm]current page.north west)
      {{\fontsize{48}{48}\selectfont\textsc{Atlas}}};
    \node[anchor=north west,text=AtlasIvory,inner sep=0pt]
      at ([xshift=17mm,yshift=-45mm]current page.north west)
      {{\fontsize{29}{32}\selectfont\textsc{of Pursuit}}};
    \node[anchor=north west,text=AtlasGoldLight,font=\normalsize\itshape,
      text width=68mm,align=left,inner sep=0pt]
      at ([xshift=17mm,yshift=-59mm]current page.north west)
      {Of trees and territories;\\of retreat, and inevitable capture.};
    \node at ([xshift=121mm,yshift=-44mm]current page.north west)
      {\AtlasGrandAstrolabe[57mm]};
    \draw[AtlasGoldLight,line width=.35pt]
      ([xshift=96mm,yshift=-19mm]current page.north west)--
      ([xshift=96mm,yshift=-71mm]current page.north west);
    \node[anchor=south west,text=AtlasGoldLight,font=\scriptsize]
      at ([xshift=17mm,yshift=11mm]current page.south west)
      {\textls[120]{\textsc{The mathematical cabinet\quad •\quad MMXXVI}}};
    \node[anchor=south east] at ([xshift=-14mm,yshift=9mm]current page.south east)
      {\pgfornament[width=23mm,color=AtlasGoldLight]{71}};
  \end{tikzpicture}
\end{frame}}

\section{Prolegomena}
\begin{frame}{Two landscapes of pursuit}
  \framesubtitle{An engraved frontispiece to the argument}
  \begin{columns}[c,onlytextwidth]
    \begin{column}{.53\textwidth}
      \centering\resizebox{.90\linewidth}{!}{\GrandTree{1}}\par
      \smallskip\AtlasGrandRubric{I. The tree}\par
      {\small\itshape Every route has a single passage.}
    \end{column}
    \begin{column}{.41\textwidth}
      \centering\resizebox{.78\linewidth}{!}{\GrandCycle{180}{0}{-1}}\par
      \AtlasGrandRubric{II. The circuit}\par
      {\small\itshape A road returns upon itself.}
    \end{column}
  \end{columns}
\end{frame}

\begin{frame}{The Argument}
  \begin{columns}[T,onlytextwidth]
    \begin{column}{.12\textwidth}\AtlasGrandInitial{A}\end{column}
    \begin{column}{.83\textwidth}
      \large pursuer walks upon a finite graph. The fugitive sees every
      move and chooses a reply. What property of the landscape makes
      capture inevitable?
    \end{column}
  \end{columns}
  \medskip
  \begin{tabular}{@{}p{10mm}>{\raggedright\arraybackslash}p{35mm}>{\raggedright\arraybackslash}p{74mm}@{}}
    \toprule
    \textcolor{AtlasRubric}{I} & \textsc{The field} & Rules, positions, and the question.\\[2mm]
    \textcolor{AtlasRubric}{II} & \textsc{The narrowing world} & A tree, a chase, and a decreasing integer.\\[2mm]
    \textcolor{AtlasRubric}{III} & \textsc{The returning road} & Where the proof fails; what a second cop repairs.\\
    \bottomrule
  \end{tabular}
  \medskip
  \centerline{\AtlasGrandFleuron[25mm]}
\end{frame}

\section{The field}
\AtlasGrandChapter{I}{The field}{First establish the rules of the encounter.}

\begin{frame}{The terms of the encounter}
  \begin{columns}[T,onlytextwidth]
    \begin{column}{.62\textwidth}
      \begin{enumerate}[<+->]
        \item The cop chooses a vertex. The robber chooses next.
        \medskip
        \item They alternate turns, with the cop moving first.
        \medskip
        \item A move follows one edge, or stays at the same vertex.
        \medskip
        \item Capture occurs when they occupy the same vertex.
      \end{enumerate}
    \end{column}
    \begin{column}{.32\textwidth}
      \AtlasGrandScholium{Scholium}{Both players see the whole graph and
        each other's positions. This is a question of strategy.}
      \centering\medskip\AtlasGrandFleuron[20mm]
    \end{column}
  \end{columns}
\end{frame}

\begin{frame}{A question fit for a single line}
  \centering
  \AtlasGrandFleuron[30mm]\par\bigskip
  {\fontsize{25}{31}\selectfont\itshape
    When does one pursuer suffice?\par}
  \bigskip
  {\color{AtlasRubric}\fontsize{30}{34}\selectfont $c(G)=1\;?$\par}
  \bigskip
  {\small\scshape The cop number is the least number guaranteeing capture.}
\end{frame}

\section{The narrowing world}
\AtlasGrandChapter{II}{The narrowing world}{Every advance leaves less room for retreat.}

\begin{frame}[label=grand-trees]{Propositio I\enspace—\enspace Of finite trees}
  \begin{theorem}[The tree principle]\label{grand:tree}
    \Large One cop can catch a robber on any finite tree.
  \end{theorem}
  \medskip
  \begin{columns}[T,onlytextwidth]
    \begin{column}{.63\textwidth}
      \AtlasGrandKicker{The construction}
      Follow the unique path toward the robber. The rest of the proof
      is an account of what the robber can never recover.
    \end{column}
    \begin{column}{.31\textwidth}
      \AtlasGrandScholium{In the margin}{Finiteness will end the chase.
      Uniqueness of paths will make progress irreversible.}
    \end{column}
  \end{columns}
\end{frame}

\begin{frame}{Tabula I\enspace—\enspace The chase}
  \framesubtitle{The plate stays still; the territory contracts}
  \begin{columns}[T,onlytextwidth]
    \begin{column}{.55\textwidth}
      \begin{overlayarea}{\linewidth}{5.4cm}
        \centering
        \only<1|handout:0>{\resizebox{.88\linewidth}{!}{\GrandTree{1}}}%
        \only<2|handout:0>{\resizebox{.88\linewidth}{!}{\GrandTree{2}}}%
        \only<3|handout:0>{\resizebox{.88\linewidth}{!}{\GrandTree{3}}}%
        \only<4|handout:1>{\resizebox{.88\linewidth}{!}{\GrandTree{4}}}%
        \only<5|handout:0>{\resizebox{.88\linewidth}{!}{\GrandTree{5}}}%
        \only<6|handout:0>{\resizebox{.88\linewidth}{!}{\GrandTree{6}}}%
        \par\smallskip
        {\scriptsize C: the pursuer\quad R: the fugitive\par
          Gold rings: the remaining region. Dashed roads: lost territory.}
      \end{overlayarea}
    \end{column}
    \begin{column}{.39\textwidth}
      \begin{overlayarea}{\linewidth}{5.4cm}
        \small\raggedright
        \only<1|handout:0>{\GrandBeat{I. The whole field}
          {The cop stands at $a$, the robber at $d$.
          Seven vertices remain reachable without crossing the cop.}{7}}
        \only<2|handout:0>{\GrandBeat{II. The first advance}
          {The cop takes $b$. The roads to $a$, $f$, and $g$
          are closed to the robber.}{4}}
        \only<3|handout:0>{\GrandBeat{III. The retreat}
          {The robber moves from $d$ to $e$.
          Retreat changes a position, but restores no territory.}{4}}
        \only<4|handout:0>{\GrandBeat{IV. The narrowing}
          {The cop takes $c$; the robber waits at $e$.
          The road to $h$ is lost as well.}{2}}
        \only<5|handout:0>{\GrandBeat{V. The last refuge}
          {The cop takes $d$. At $e$, the robber's only safe reply is to wait.}{1}}
        \only<6|handout:0>{\GrandBeat{VI. The meeting}
          {The cop takes $e$. The two pieces coincide.
          The remaining region is empty: capture.}{0}}
        \only<beamer:0|handout:1>{%
          \AtlasGrandKicker{A record of the chase}
          \begin{tabular}{@{}crrc@{}}\toprule
            Step & C & R & Region\\\midrule
            I & $a$ & $d$ & 7\\
            II & $b$ & $d$ & 4\\
            III & $b$ & $e$ & 4\\
            IV & $c$ & $e$ & 2\\
            V & $d$ & $e$ & 1\\
            VI & $e$ & $e$ & 0\\\bottomrule
          \end{tabular}\par\smallskip
          Step IV is pictured. The robber waits in IV and V.

        }
      \end{overlayarea}
    \end{column}
  \end{columns}
\end{frame}

\begin{frame}{Demonstratio\enspace—\enspace An irreversible loss}
  \begin{columns}[T,onlytextwidth]
    \begin{column}{.53\textwidth}
      \begin{overlayarea}{\linewidth}{5.8cm}
        \only<1|handout:0>{\resizebox{.90\linewidth}{!}{\GrandTree{3}}}%
        \only<2-|handout:1>{\resizebox{.90\linewidth}{!}{\GrandTree{4}}}%
        \par\smallskip
        \centering\small
        \only<1|handout:0>{$S=\{c,d,e,h\}$}
        \only<2-|handout:1>{$S'=\{d,e\}$}
      \end{overlayarea}
    \end{column}
    \begin{column}{.42\textwidth}
      \begin{overlayarea}{\linewidth}{5.8cm}
        \small\raggedright
        \AtlasGrandKicker{I. The region}
        Delete the cop's vertex. Let $S$ be the robber's component.
        \par\medskip
        \uncover<2->{%
          \AtlasGrandKicker{II. Its boundary}
          Step to $v$ on the unique path. Every safe reply leaves
          \par\smallskip
          {\color{AtlasRubric}\large $S'\subseteq S\setminus\{v\}.$\par}
          \medskip
        }
        \uncover<3->{%
          \AtlasGrandKicker{III. Finiteness}
          Since $v\in S$, we have $|S'|<|S|$.
          A positive integer cannot fall forever.
          \par\smallskip\hfill\usebeamertemplate{qed symbol}
        }
      \end{overlayarea}
    \end{column}
  \end{columns}
\end{frame}

\begin{frame}{Tabula II\enspace—\enspace The account of territory}
  \centering
  \resizebox{.80\linewidth}{!}{\GrandDescent}\par
  \smallskip
  \begin{minipage}{.85\linewidth}
    \AtlasGrandScholium{What the record counts}{The bars record the region
    after successive cop advances in this chase. A robber move can leave
    the count unchanged; it cannot make the lost region return.}
  \end{minipage}
\end{frame}

{\AtlasGrandNight
\begin{frame}[plain,noframenumbering]
  \centering
  \pgfornament[width=45mm,color=AtlasGoldLight]{71}\par\bigskip
  {\color{AtlasIvory}\fontsize{27}{34}\selectfont\itshape
    The proof is a landscape\\
    that grows smaller\\
    with every advance.\par}
  \bigskip
  {\color{AtlasGoldLight}\Large $7\;>\;4\;>\;2\;>\;1\;>\;0$\par}
  \bigskip
  {\color{AtlasGoldLight}\scriptsize\textls[160]{\textsc{An interlude on descent}}\par}
\end{frame}}

\section{The returning road}
\AtlasGrandChapter{III}{The returning road}{A single circuit changes the nature of retreat.}

\begin{frame}{Tabula III\enspace—\enspace A road back}
  \begin{columns}[c,onlytextwidth]
    \begin{column}{.48\textwidth}
      \begin{overlayarea}{\linewidth}{5.8cm}
        \centering
        \only<1|handout:0>{\GrandCycle{180}{0}{-1}}%
        \only<2|handout:0>{\GrandCycle{120}{0}{-1}}%
        \only<3|handout:1>{\GrandCycle{120}{300}{-1}}%
      \end{overlayarea}
    \end{column}
    \begin{column}{.46\textwidth}
      \begin{overlayarea}{\linewidth}{5.8cm}
        \small\raggedright
        \AtlasGrandKicker{The missing invariant}
        Removing one vertex leaves a connected path of five vertices.
        The region's size does not change.
        \par\smallskip
        \uncover<2->{The cop moves one edge closer.\par\smallskip}
        \uncover<3->{The robber moves away and restores the opposite position.\par}
        \smallskip
        \AtlasGrandScholium{The distinction}{Distance can fall for a turn
        without territory being lost.}
      \end{overlayarea}
    \end{column}
  \end{columns}
\end{frame}

\begin{frame}{Tabula IV\enspace—\enspace Two boundaries}
  \begin{columns}[c,onlytextwidth]
    \begin{column}{.48\textwidth}
      \begin{overlayarea}{\linewidth}{5.8cm}
        \centering
        \only<1|handout:0>{\GrandCycle{180}{60}{0}}%
        \only<2|handout:1>{\GrandCycle{120}{60}{0}}%
        \only<3|handout:0>{\GrandCycle{60}{60}{0}}%
      \end{overlayarea}
    \end{column}
    \begin{column}{.46\textwidth}
      \begin{overlayarea}{\linewidth}{5.8cm}
        \small\raggedright
        \AtlasGrandKicker{Hold; then advance}
        One cop holds $0^\circ$. The other advances along the robber's path.
        \par\smallskip
        \uncover<2->{At $120^\circ$, the moving cop leaves only a wait at $60^\circ$.\par\smallskip}
        \uncover<3->{The next advance captures the robber.\par}
        \smallskip
        \AtlasGrandScholium{The repaired argument}{The fixed cop cuts the cycle.
        The moving cop shortens the path.}
      \end{overlayarea}
    \end{column}
  \end{columns}
\end{frame}

\begin{frame}{The proposition, the example, the caution}
  \begin{columns}[T,onlytextwidth]
    \begin{column}{.58\textwidth}
      \begin{definition}[Cop number]
        $c(G)$ is the least number of cops that can guarantee capture on $G$.
      \end{definition}
      \begin{exampleblock}{The tree}
        Proposition~\ref{grand:tree} gives $c(T)=1$ for every finite tree.
      \end{exampleblock}
    \end{column}
    \begin{column}{.36\textwidth}
      \begin{alertblock}{The circuit}
        For a cycle of length at least four, $c(C_n)=2$.
      \end{alertblock}
      \centering\medskip\AtlasGrandFleuron[21mm]\par\medskip
      {\small\itshape A hypothesis is part of the theorem's shape.}
    \end{column}
  \end{columns}
\end{frame}

\section{The cabinet}
\begin{frame}{A leaf for the marginal reader}
  \begin{columns}[T,onlytextwidth]
    \begin{column}{.65\textwidth}
      \begin{columns}[T,onlytextwidth]
        \begin{column}{.2\textwidth}\AtlasGrandInitial{E}\end{column}
        \begin{column}{.75\textwidth}
          very proof has a measure. Here it is the number of vertices
          left in a territory that cannot be regained.
        \end{column}
      \end{columns}
      \bigskip
      {\large\itshape The drawing shows what is lost.\\
        The inequality shows why it matters.\par}
      \bigskip
      {\color{AtlasRubric}\LARGE $S'\subsetneq S$\par}
      \medskip\AtlasGrandFleuron[28mm]
    \end{column}
    \begin{column}{.29\textwidth}
      \AtlasGrandScholium{Scholium I}{The region is a component, not a ball
      of a prescribed radius.}
      \AtlasGrandScholium{Scholium II}{The chase illustrates a strategy.
      The invariant must cover every reply.}
    \end{column}
  \end{columns}
\end{frame}

\begin{frame}{Colophon\enspace—\enspace What this edition attempts}
  \begin{columns}[T,onlytextwidth]
    \begin{column}{.6\textwidth}
      \AtlasGrandKicker{The type}
      Garamond text and mathematics; rubric, small capitals,
      illuminated letters, and Roman folios.
      \par\medskip
      \AtlasGrandKicker{The plate}
      Compass rings, armillary linework, engraved roads, medallion pieces,
      hatched quantities, and a dial that runs down to capture.
      \par\medskip
      \AtlasGrandKicker{The sheet}
      Laid lines, chain lines, gilded borders, corner ornaments,
      and a dark reverse for the chapter openings.
    \end{column}
    \begin{column}{.34\textwidth}
      \AtlasGrandScholium{The experiment}{How much ceremony can a mathematical
      presentation carry before the apparatus begins to rival the argument?}
      \centering\AtlasGrandFleuron[24mm]\par\medskip
      {\scriptsize All illustrations are vector constructions.\\
      The movements are speaker-paced overlays.}
    \end{column}
  \end{columns}
\end{frame}

{\AtlasGrandNight
\begin{frame}[plain,noframenumbering]
  \centering
  \AtlasGrandAstrolabe[32mm]\par\medskip
  {\color{AtlasIvory}\fontsize{38}{40}\selectfont\textsc{Finis}\par}
  \medskip
  {\color{AtlasGoldLight}\large\itshape The circuit invites another argument.\par}
  \bigskip
  \pgfornament[width=38mm,color=AtlasGoldLight]{71}\par\smallskip
  {\color{AtlasGoldLight}\scriptsize\textls[160]{\textsc{Set for the screen\quad •\quad MMXXVI}}\par}
\end{frame}}

\end{document}
